\documentclass[11pt]{article}
\usepackage[margin=1in]{geometry}
\usepackage{amsmath,amssymb,amsthm}
\usepackage{booktabs}
\usepackage{listings}
\usepackage{xcolor}
\usepackage[numbers]{natbib}
\usepackage{hyperref}

\newtheorem{theorem}{Theorem}
\newtheorem{conjecture}{Conjecture}
\theoremstyle{remark}
\newtheorem{remark}{Remark}

\lstset{basicstyle=\ttfamily\small,columns=fullflexible,breaklines=true,
  frame=single,keywordstyle=\color{blue!60!black},commentstyle=\color{gray},
  morekeywords={def,theorem,lemma,namespace,end,example,by,fun,import,noncomputable,section,open,intro,have,rw,at,exact,axiom},
  morecomment=[l]{--},
  literate={ℝ}{{$\mathbb{R}$}}1 {→}{{$\to$}}1 {∀}{{$\forall$}}1 {∃}{{$\exists$}}1
           {¬}{{$\neg$}}1 {≤}{{$\le$}}1 {≥}{{$\ge$}}1 {↔}{{$\leftrightarrow$}}1
           {∧}{{$\wedge$}}1 {δ}{{$\delta$}}1 {·}{{$\cdot$}}1 {−}{{$-$}}1 {⁻}{{$^{-}$}}1}

\title{A Correction Note on an Early Intelligence Bound Source:\\
The Unrestricted Axiom Is False as Written}
\author{Justin Hart\\Viridis LLC}
\date{14 November 2026}

\begin{document}
\maketitle

\begin{abstract}
An early Viridis Lean source of January 2026 declares an axiom, \texttt{intelligence\_bound},
stating $\dot I\le \frac{P}{k_B T\ln 2}\,D$ for every real number $\dot I$, under the guards
$k_B>0$, $T>0$, $P\ge0$ and $0\le D\le1$, with no hypothesis linking $\dot I$ to any system.
The instance $\dot I=1$, $P=D=0$, $k_B=T=1$ meets every guard and reduces the conclusion to
$1\le0$. We give a Lean~4 proof, against the sealed statement
\texttt{P5.UnrestrictedBoundFalse}, that the universally quantified proposition is false, and a
companion Lean lemma showing that, taken as a hypothesis, it proves every proposition. An
environment that assumes the axiom is therefore inconsistent. According to a reading of the
corpus, no declaration in the reviewed sources depended on it, so no recorded result is affected.
We describe what a well-posed successor needs: a predicate that ties the rate to a modelled
physical system, and separate empirical validation. This note is about one piece of formal
text. It says nothing about the Intelligence Bound as a physical conjecture. A second correction,
on finite compensation of a loss of $D$, is restated here as a conjecture pending its own proof.
\end{abstract}

\section{Introduction}
A machine-checked development is only as reliable as the match between its formal text and what
is said about it. A Lean \texttt{axiom} is not checked at all: the kernel accepts it as given. If
the axiom is false, every theorem proved in its presence is vacuous, because a false hypothesis
proves anything. A review of the Viridis research corpus on 11~November~2026 found two early Lean
declarations whose formal text diverges from their surrounding prose. Project P5 checks both with
the Lean kernel. The first session (P5-T001) proposed exact statements, and these were then sealed.
This note covers the first finding: the historical \texttt{intelligence\_bound} axiom is false
as written.

The question matters for record keeping. The axiom was presented as a formal version of a
physical bound on information processing. As written, though, it is a claim about all real
numbers, and it fails for elementary reasons. Correcting it is not new mathematics. The aim is to
keep false text from being cited as a foundation, and to state what a correct successor would
have to contain.

\section{Background and related work}
Landauer's principle~\citep{landauer1961} associates a minimal dissipation of $k_B T\ln 2$ with
each logically irreversible bit operation at temperature $T$. Dividing a power budget $P$ by this
energy gives an upper bound on the rate of irreversible bit operations. The early sources define,
over the reals,
\[
  \mathrm{landauerEnergy}(k_B,T)=k_B T\ln 2,\qquad
  \mathrm{maxBitRate}(P,k_B,T)=\frac{P}{k_B T\ln 2},
\]
and multiply the result by a diversity factor $D\in[0,1]$. We use Lean~4~\citep{lean4} and
Mathlib~\citep{mathlib2020}. The historical file \texttt{Gaia\_Intelligence\_Theorem.lean} is
identified by its SHA-256 prefix \texttt{37c017145dcf\ldots2098c4}. It is quoted, not imported:
the sealed statement restates the two definitions over Mathlib alone.

\section{The historical axiom}
The project README quotes the declaration as follows:
\begin{lstlisting}
axiom intelligence_bound :
  ∀ (dI_dt P D k_B T : ℝ),
  k_B > 0 → T > 0 → P ≥ 0 → D ≥ 0 → D ≤ 1 →
  dI_dt ≤ max_bit_rate P k_B T * D
\end{lstlisting}
The rate \texttt{dI\_dt} is a bound variable that ranges over every real number. Nothing in the
statement says whose rate it is. The axiom therefore asserts that every real number lies below
$\mathrm{maxBitRate}(P,k_B,T)\cdot D$, which cannot hold, because the right-hand side is a fixed
real number once the other variables are fixed.

\section{Result}
\begin{theorem}[sealed statement \texttt{P5.UnrestrictedBoundFalse}]\label{thm:main}
It is not the case that for all real $\dot I,P,D,k_B,T$ with $k_B>0$, $T>0$, $P\ge0$, $D\ge0$
and $D\le1$ we have $\dot I\le\mathrm{maxBitRate}(P,k_B,T)\cdot D$.
\end{theorem}
\begin{proof}[Counterinstance]
Take $\dot I=1$, $P=0$, $D=0$, $k_B=T=1$. Every guard holds. The right-hand side is
$\frac{0}{\ln 2}\cdot 0=0$, so the conclusion would be $1\le0$, which is false.
\end{proof}

We give a Lean proof of Theorem~\ref{thm:main}. The automated checker run during the session
passed all five stages, and the only axioms used are \texttt{propext}, \texttt{Classical.choice}
and \texttt{Quot.sound}. The candidate file is:
\begin{lstlisting}
lemma bound_zero_power (D kB T : ℝ) : maxBitRate 0 kB T * D = 0 := by
  simp [maxBitRate]

theorem main : ViridisStatements.P5.UnrestrictedBoundFalse.Statement := by
  intro h
  -- Counterinstance: dIdt = 1, P = 0, D = 0, kB = T = 1.
  have h1 := h 1 0 0 1 1 one_pos one_pos le_rfl le_rfl zero_le_one
  rw [bound_zero_power] at h1
  exact absurd h1 (by norm_num)

theorem bound_proves_anything
    (h : ∀ dIdt P D kB T : ℝ, kB > 0 → T > 0 → P ≥ 0 → D ≥ 0 → D ≤ 1 →
      dIdt ≤ maxBitRate P kB T * D)
    (p : Prop) : p :=
  absurd h main
\end{lstlisting}

\begin{remark}[The witness is not degenerate]
The witness uses $P=0$ and $D=0$, but the falsity does not depend on these boundary values. For
any admissible $P,D,k_B,T$, the rate $\dot I=\mathrm{maxBitRate}(P,k_B,T)\cdot D+1$ violates the
bound. In particular this works for $P>0$ and $D=1$. The witness also does not rely on Lean's
convention $x/0=0$: here the denominator is $\ln 2\neq0$.
\end{remark}

\paragraph{Consequence: inconsistency.} The lemma \texttt{bound\_proves\_anything} makes the
consequence explicit. If the universally quantified bound is available as a hypothesis, or as an
axiom, any proposition $p$ follows, including $p$ and $\neg p$ alike (\emph{ex falso
quodlibet}). A Lean environment that declares \texttt{intelligence\_bound} as an axiom is
therefore inconsistent. Every statement can be proved in it, and a successful compilation in that
environment certifies nothing.

\paragraph{Impact on the corpus.} The 2026-11-11 review reports that no declaration in
\texttt{Gaia\_Intelligence\_Theorem.lean} uses the axiom and that no other file in the reviewed
corpus imports it. If that is right, no recorded theorem inherits the inconsistency. This rests on
reading the sources. It is not established by the Lean proof above, which concerns only the
restated proposition.

\section{What a well-posed successor needs}
The negative control in the sealed module shows what is missing. Once the rate is tied to the
bound by hypothesis, the implication holds trivially. A meaningful successor therefore cannot
quantify the rate freely. It needs:
\begin{enumerate}
\item \emph{A predicate tying the rate to a physical system.} For example, a structure modelling
  a system with power $P$ at temperature $T$, whose information rate is \emph{defined} from that
  model, perhaps as a count of irreversible bit operations per unit time. The bound then becomes a
  theorem about that structure, derived from an explicit Landauer-type hypothesis, rather than an
  axiom about arbitrary reals. The role of $D$ also has to be defined inside the model, not just
  multiplied in.
\item \emph{Separate empirical validation.} Whether such a model describes real agents or
  ecosystems is an empirical question. No Lean proof answers it. Its premises must be stated and
  tested against data, and they must be labelled as assumptions until that is done.
\end{enumerate}
Neither item is attempted here.

\section{The second correction (pending)}
For the record, the second P5 finding is restated without change. Let $\mathrm{cap}(P,D,C)=PD/C$.
The historical theorem \texttt{no\_finite\_compensation} actually proves that the finite power
$PD/(D-\Delta D)$ restores the capacity after a partial loss $0<\Delta D<D$.
\begin{conjecture}[sealed statement \texttt{P5.CompensationUnderCeiling}]
Let $C>0$ and $P>0$. (1) If $0<\Delta D<D$, then $P'=PD/(D-\Delta D)$ is the unique power with
$\mathrm{cap}(P',D-\Delta D,C)=\mathrm{cap}(P,D,C)$. (2) If moreover $P\le P_{\max}$, then such a
$P'\le P_{\max}$ exists if and only if $\Delta D\le D(1-P/P_{\max})$. (3) For $D>0$,
$PD/(D-\Delta D)\to\infty$ as $\Delta D\to D^-$. (4) For $D>0$ and every $P'$,
$\mathrm{cap}(P',0,C)=0<\mathrm{cap}(P,D,C)$.
\end{conjecture}
\noindent The proof is task P5-T003. No proof of it is claimed in this note.

\section{Limitations}
This is a correction of one piece of formal text. It is elementary, and it is not a result about
the Intelligence Bound conjecture. Showing that the unrestricted text is false says nothing about
whether a correctly posed bound holds for any physical system. It also says nothing about
whether any agent would preserve biodiversity. The statement that no corpus declaration depended
on the axiom comes from reading the sources and is not machine-checked. The sealed statement
restates the historical definitions instead of importing the historical file. Its fidelity to the
original therefore rests on the quoted text and the file hash.

\section{Reproducibility}
The proof is \texttt{proofs/P5.UnrestrictedBoundFalse.lean}. It imports only Mathlib and the
sealed module \texttt{ViridisStatements.P5.UnrestrictedBoundFalse} (SHA-256
\path{3b83eb37905660e3db1cc68c54e476f4ac5c450b245cec4675195f9cc2e0a95d}). It compiles with
toolchain \texttt{leanprover/lean4:v4.28.0} and Mathlib commit
\path{8f9d9cff6bd728b17a24e163c9402775d9e6a365}, and it can be checked with
\texttt{viridis-lean-check proofs/P5.UnrestrictedBoundFalse.lean P5.UnrestrictedBoundFalse}.
No data or external code is used.

\section{Changelog}
\begin{itemize}
\item 2026-11-13 (P5-T001): first version. Proposed both statements with sanity checks and
negative controls, and noted the conditions the negative controls need.
\item 2026-11-14 (P5-T002): refocused the note on the first correction. Added a Lean proof of
the sealed \texttt{P5.UnrestrictedBoundFalse} and an explosion lemma, explained the
inconsistency and its impact on the corpus, and set out the requirements for a well-posed
successor. The second correction remains a conjecture.
\end{itemize}

\bibliographystyle{plainnat}
\bibliography{refs}
\end{document}
